% macro conflict, first definition kept: \newcommand{\Homext}{\mathrm{Hom}\,\mathrm{ext}}
\documentclass[11pt]{article}
\usepackage{fontspec}
\usepackage[french]{babel}
\usepackage{amsmath}
\usepackage{amssymb}
\usepackage{amsthm}
\usepackage{mathrsfs}
\usepackage{tikz-cd}
\usepackage{enumitem}

\newcommand{\printedpage}[1]{\marginpar{\footnotesize p.~#1}}
\newcommand{\msp}[1]{\marginpar{\footnotesize [#1]}}
\newcommand{\Grp}{\underline{\mathrm{Grp}}}
\newcommand{\uHom}{\underline{\mathrm{Hom}}}
\newcommand{\uGrp}{\underline{\mathrm{Grp}}}
\newcommand{\uuGrp}{\underline{\underline{\mathrm{Grp}}}}
\newcommand{\fl}{\longrightarrow}
\newcommand{\Ch}{\mathrm{Ch}}
\newcommand{\Hot}{\mathrm{Hot}}
\newcommand{\Hom}{\mathrm{Hom}}
\newcommand{\HotO}{\mathrm{Hot}_0^{\wedge}}
\newcommand{\MOP}{\mathcal{M}_{0\,\Pi}}
\newcommand{\M}{\mathcal{M}}
\newcommand{\uCh}{\underline{\mathrm{Ch}}}
\newcommand{\Homext}{\mathrm{Hom\,ext}}
\newcommand{\Lotimes}{\overset{\mathbf{L}}{\otimes}}
\newcommand{\cercle}[1]{\tikz[baseline=(c.base)]\node[draw,circle,inner sep=1pt](c){#1};}
\newcommand{\Chu}{\underline{\mathrm{Ch}}}
\newcommand{\D}{\mathbf{D}}
\newcommand{\Cat}{\mathrm{Cat}}
\newcommand{\Fl}{\operatorname{Fl}}
\newcommand{\Ob}{\operatorname{Ob}}
\newcommand{\bDelta}{\boldsymbol{\Delta}}
\begin{document}
% ---- from Der-ChVIII-567_p001-005.tex
% ===================== page 1 =====================
\printedpage{}
% head: (aucun ; page de titre non numérotée)

\begin{center}
{\Large\textsf{LES DÉRIVATEURS}}

\bigskip

\textbf{ALEXANDRE GROTHENDIECK}

\bigskip\bigskip

\textbf{Chapitre VIII}

\medskip

\textbf{1-types d'homotopie relatifs : leur intégration \dots}
\end{center}

\clearpage

% ===================== page 2 =====================
\printedpage{}
% head: (aucun ; page non numérotée)

Ce texte a été déchiffré et transcrit en \LaTeXe{} par M. Künzer. Il a été édité par M. Künzer, J. Malgoire, et G. Maltsiniotis. La transcription est aussi fidèle que possible au manuscrit. Pour les quelques corrections évidentes, ou rares commentaires des éditeurs, ainsi que pour la numérotation originale des pages du manuscrit, les caractères de machine à écrire \texttt{[typewriter]} entre crochets sont utilisés. Un point d'interrogation entre crochets signifie que l'on n'est pas sûr du mot qui précède.

Cette édition est provisoire. Les remarques, commentaires et corrections sont bienvenus. Envoyer un message à :

\textbf{[e-mail address removed]}

\hfill G. Maltsiniotis\hspace*{2cm}

\clearpage

% ===================== page 3 =====================
\printedpage{1}
% head: A. Grothendieck : Dérivateurs, Ch. VIII, Édition des Universités de Montpellier II et Paris VII
% foot: Version 2 février 2002

\msp{page 1}

\section*{1\quad Intégration des $1$-types d'homotopie relatifs sur une catégorie.}

Les types d'homotopie relatifs sur une (petite) catégorie $S$ sont définis par les \textit{champs} (en catégories) sur $S^{\wedge}$, qui correspondent essentiellement aux contrafoncteurs $S^o \longrightarrow (\mathrm{Cat})$, ou encore aux catégories $X$ sur $S$ qui sont Cat-fibrées sur $S$. (Dans $(\mathrm{Cat}/S)_{\mathrm{fib}}$, on inverse les flèches qui sont des Hot-équivalences sur les fibres.) D'autre part, les 1-types d'homotopie relatifs sur $S$ doivent se décrire par des \textit{champs en groupoïdes} sur $S$. Si on a un type d'homotopie relatif, défini par
\[ X \longrightarrow S, \qquad\qquad \text{catégorie Cat-fibrée sur } S, \]
ou par
\[
\begin{aligned}
S^o \quad &\longrightarrow \quad \mathrm{Cat}\\
s \quad &\longmapsto \quad X_s\ ,
\end{aligned}
\]
le 1-type d'homotopie associé est défini par
\[ s \longmapsto \Pi(X_s), \qquad\qquad S^o \longrightarrow (\mathrm{Grp}), \qquad\qquad \text{catégorie des groupoïdes,} \]
où $\Pi(X_s)$ est le groupoïde fondamental de $X_s$. On note la catégorie fibrée correspondante sur $S$ par \fbox{$\Pi(X/S)$}. Il est vrai que si
\[ X \ \longrightarrow\ X' \]
est un homomorphisme de catégories fibrées sur $S$ qui induise une Hot-équivalence sur les fibres, elle induit une équivalence fibrée
\[ \Pi(X/S) \ \longrightarrow\ \Pi(X'/S), \]
mais non un isomorphisme. Il y aurait donc du travail à faire pour définir un foncteur
\[
\underbrace{\mathrm{Hot}(S)}_{\text{\footnotesize\begin{tabular}{c}catégorie des\\ types d'homotopie\\ relatifs sur $S$\end{tabular}}}
\quad\xrightarrow{\ \xi\,\longmapsto\,\Pi(\xi/S)\ }\quad
\underbrace{\Grp(S)}_{\text{\footnotesize\begin{tabular}{c}catégorie des\\ groupoïdes\\ fibrés sur $S$\\ (ou catégories\\ des champs en\\ groupoïdes\\ sur $S^{\wedge}$)\end{tabular}}}
\quad .
\]

Mais ce n'est pas pour l'instant mon propos. Plutôt, de donner une formule de transitivité en homologie non commutative, analogue à $R(gf)_* \simeq Rg_*Rf_*$ en cohomologie.

\clearpage

% ===================== page 4 =====================
\printedpage{2}
% head: A. Grothendieck : Dérivateurs, Ch. VIII, Édition des Universités de Montpellier II et Paris VII
% foot: Version 2 février 2002

\msp{page 2}

Si
\[ f : S \longrightarrow T \]
est un foncteur, on va définir un foncteur
\[ f_! :\ \Grp(S) \ \longrightarrow\ \Grp(T) \]
(image directe de 1-types d'homotopie relatifs). Mais notons d'abord que si
\[ p :\ X \ \longrightarrow\ S \]
est une $S$-catégorie \textit{quelconque} (pas nécessairement Cat-fibrée), on lui associe néanmoins un type d'homotopie relatif, via la catégorie fibrée $\tilde{X}^{(S)}$, ou $\tilde{X}$, ayant comme fibres les $s\backslash X$ (en notant que $s' \longrightarrow s$ dans $S$ définit bien $s\backslash X \longrightarrow s'\backslash X$) (\footnote{Ayant $\tilde{X}$, on définit $\Pi(X/S) \overset{\text{déf}}{=} \Pi(\tilde{X}/S)$.}). On a un $S$-foncteur canonique
\[ X \xrightarrow{\ i\ } \tilde{X}, \qquad\qquad x \longmapsto (x, p(x), \underbrace{p(x)}_{=\,\text{‘}\,s\,\text{’}} \xrightarrow{\ \mathrm{id}\ } p(x)). \]
C'est une $S$-Hot-équivalence si $X$ était déjà fibrée. En fait, sur chaque fibre c'est un homotopisme car on définit alors $\tilde{X}_{s_0} \xrightarrow{\ j\ } X_{s_0}$ de façon évidente, de façon que
\[
X_{s_0} \underset{j}{\overset{i}{\rightleftarrows}} \tilde{X}_{s_0}, \qquad ji \simeq \mathrm{id}, \qquad ij \underbrace{\xrightarrow{\text{homotopie simple}}}_{\text{\footnotesize\begin{tabular}{l}NB cette flèche\\ n'est pas \texttt{[un]}\\ isomorphisme\\ en général.\end{tabular}}} \mathrm{id}.
\]
Revenons alors à $f : S \longrightarrow T$, et considérons un groupoïde fibré sur $S$,
\[ G \ \longrightarrow\ S, \]
ce n'est pas une catégorie fibrée sur $T$, mais on peut lui associer
\[ t \longmapsto t\backslash G, \qquad\qquad T^o \longrightarrow (\mathrm{Cat}), \]
d'où une catégorie fibrée sur $T$, dont on prend le $\Pi$ relatif. Mais il est plus simple de prendre directement
\[ t \ \longmapsto\ \Pi(t\backslash G), \]
qui définit une catégorie fibrée en groupoïdes sur $T$.

\clearpage

% ===================== page 5 =====================
\printedpage{3}
% head: A. Grothendieck : Dérivateurs, Ch. VIII, Édition des Universités de Montpellier II et Paris VII
% foot: Version 2 février 2002

\msp{page 3}

Ainsi on trouve bien un foncteur (en fait, même un 2-foncteur de 2-catégories)
\[ f_! :\ \Grp(S) \ \longrightarrow\ \Grp(T). \]
Si
\[ u :\ G \ \longrightarrow\ G' \]
est un homomorphisme de groupoïdes fibrés sur $S$, qui induise des Cat-équivalences sur les fibres, alors
\[ f_!(u) :\ f_!(G) \ \longrightarrow\ f_!(G') \]
est aussi une Cat-équivalence sur les fibres. (Cela est formel, à partir du fait que $f_!$ est un 2-foncteur, compte tenu que si on a une équivalence sur les fibres, alors il y a un quasi-inverse -- c'est une histoire générale de catégories fibrées \dots)

Soit maintenant
\[ p :\ X \ \longrightarrow\ S, \]
et prenons le cas $G = \Pi(X/S)$. Je dis qu'on a alors une $T$-équivalence
\[ \boxed{\ f_!(\Pi(X/S)) \ \simeq\ \Pi(X/T),\ } \tag*{$(*)$} \]
en d'autres termes, pour connaître $\Pi(X/T)$ (à $T$-équivalence près), il suffit de connaître $\Pi(X/S)$, et de lui appliquer le foncteur $f_!$ d'intégration des 1-types d'homotopie relatifs, pour $f : S \longrightarrow T$.

Notons que l'on a trivialement
\[ \Pi(X/S) \ =\ p_!(e_X), \]
où $e_X$ désigne le $X$-groupoïde fibré trivial (toutes les fibres ponctuelles). Donc la formule précédente s'écrit comme une formule de transitivité
\[ f_!(p_!(\Gamma)) \ \simeq\ (fp)_!(\Gamma), \]
\msp{page 4}
où $\Gamma$ est le groupoïde fibré trivial sur $X$. Cela suggère de généraliser la formule $(*)$ en
\[ \boxed{\ (fp)_!(\Gamma) \ \simeq\ f_!(p_!(\Gamma))\ } \]
pour \textit{tout} groupoïde fibré sur $X$, dans une situation
\[ X \xrightarrow{\ p\ } S \xrightarrow{\ f\ } T. \]

Ce sera formel, si on établit une propriété de 2-adjonction pour le foncteur $f_!$ :

% ---- from Der-ChVIII-567_p006-010.tex
\printedpage{4}
% head: A. Grothendieck : Dérivateurs, Ch. VIII, Édition des Universités de Montpellier II et Paris VII
% foot: Version 2 février 2002

\textbf{Théorème.} \textit{Soit $f : S \fl T$ une flèche dans} (Cat)\textit{, $G$ un champ en groupoïdes sur $S$, $H$ un champ en groupoïdes sur $T$. On a alors une équivalence de catégories canonique}
\[
\boxed{\ \uHom_{\uGrp(S)}(G, f^*H) \ \simeq\ \uHom_{\uGrp(T)}(f_!G, H).\ }
\]

Les flèches d'adjonction
\[
\left\{
\begin{array}{lclcrcl}
f_!f^*H & \fl & H & \qquad\qquad & H & \in & \mathrm{Ob}\,\uGrp(T)\\[4pt]
f^*f_!G & \longleftarrow & G & \qquad\qquad & G & \in & \mathrm{Ob}\,\uGrp(S)
\end{array}
\right.
\]
devraient se définir sans problème. Pour la deuxième, il faut définir pour $s \in S$
\[
G_s \ \fl\ (f^*(f_!(G)))_s = (f_!(G))_{f(s)} \overset{\text{déf}}{=} \Pi(s\backslash G),
\]
or on a
\[
G_s \fl s\backslash G,
\]
d'où
\[
\underbrace{\Pi(G_s)}_{\begin{subarray}{l} =\ G_s\\ \text{car } G_s\\ \text{est un}\\ \text{groupoïde}\end{subarray}} \ \fl \Pi(s\backslash G),
\]
O.K. C'est fonctoriel en $s$ \dots

\msp{page 5}
Pour la flèche $f_!f^*H \fl H$ sur $T$, il faut trouver pour tout $t \in \mathrm{Ob}\,T$, et fonctoriellement en $t$,
\[
\underbrace{(f_!f^*H)_t}_{\overset{\text{déf}}{=}\Pi(t\backslash f^*(H))} \quad \fl \quad \underbrace{H_t}_{=\,\Pi(H_t)} \ ,
\]
il suffirait donc de définir
\[
t\backslash f^*(H) \ \fl\ H_t,
\]
mais on a le diagramme cartésien
\[
\begin{tikzcd}[column sep=large, row sep=large]
f^*H = H \times_S T \arrow[r] \arrow[d] & H \arrow[d] \\
S \arrow[r, "f"] & T \ni t
\end{tikzcd}
\]
\texttt{[plutôt} $H \times_T S$\texttt{]} induisant
\[
f^*H \ \fl\ H,
\]
\printedpage{5}
% head: A. Grothendieck : Dérivateurs, Ch. VIII, Édition des Universités de Montpellier II et Paris VII
% foot: Version 2 février 2002
d'où
\[
t\backslash f^*H \ \fl\ t\backslash H \qquad\quad \text{(au lieu de s'envoyer dans $H_t$),}
\]
mais d'autre part, comme $H$ est fibrée sur $T$, donc lisse sur $T$, l'inclusion
\[
H_t \fl t\backslash H
\]
est une Hot-équivalence. Ou disons plutôt : $f^*H \fl H$ est un morphisme de catégories sur $T$, induisant
\[
\underbrace{\Pi(f^*H/T)}_{\overset{\text{déf}}{=}\ f_!(f^*H)} \ \fl\ \Pi(H/T),
\]

\msp{page 6}
or le $T$-foncteur canonique
\[
H \fl \Pi(H/T) \qquad\quad (= \text{catégorie sur } T \text{ définie par } t \longmapsto \Pi(t\backslash H))
\]
est une équivalence, parce que $H$ est un fibré en groupoïdes sur $T$. Donc on peut trouver un $T$-quasi-inverse
\[
\Pi(H/T) \fl H
\]
et en composant on trouve
\[
f_!(f^*H) \ \fl\ \Pi(H/T) \ \fl\ H
\]
la flèche cherchée.

Ayant les 2-flèches d'adjonction, je m'arrête là dans les vérifications !

\printedpage{6}
% head: A. Grothendieck : Dérivateurs, Ch. VIII, Édition des Universités de Montpellier II et Paris VII
% foot: Version 2 février 2002

\section*{2\quad Questions de changement de base}

\[
\begin{tikzcd}[column sep=large, row sep=large]
S & S' \arrow[l, "q"'] \arrow[d, "f'"] \\
T \arrow[from=1-1, "f"'] & T' \arrow[l, "p"]
\arrow[from=2-2, to=1-1, bend right=50, shorten <=18pt, shorten >=18pt, "\gamma"']
\end{tikzcd}
\qquad\qquad\qquad
p^*f_!(G) \longleftarrow f'_!q^*(G)
\tag*{$(*)$}
\]

La flèche est définie, pourvu qu'on ait commutativité ou une flèche de commutativité $\gamma : pf' \fl fq$ (pas nécessairement \texttt{[un]} isomorphisme). Par adjonction, il revient au même de donner
\[
\begin{tikzcd}[row sep=large]
p_!f'_!q^*G \arrow[r] \arrow[d, dash, "\wr"', "\text{transitivité}"] & f_!(G) \\
(pf')_!q^*G \arrow[d, "\gamma"] & \\
(fq)_!q^*G \arrow[d, dash, "\text{transitivité}\ \wr"'] & \\
f_!(q_!q^*)G. \arrow[uuur, dashed, "\text{car } q_!q^*G \longrightarrow G"'] &
\end{tikzcd}
\]

On se demande quand la flèche est \texttt{[un]} isomorphisme.

\msp{page 7}
Le théorème suivant doit être vrai, et résulte de mes fourbis généraux.

% vertical bar printed in the left margin beside the following theorem
\textbf{Théorème.} \textit{Si le carré $(*)$ est} cartésien \textit{dans} (Cat)\textit{, $\gamma$ un isomorphisme, et si on suppose que $f$ soit} 1\textit{-lisse, ou $p$} 1\textit{-propre, alors le $T'$-foncteur de changement de base est une équivalence de $T'$-catégories.}

Les ‘fourbis’ généraux seraient en termes d'un dérivateur $\mathbf{D}$ ad hoc, celui des 1-types d'homotopie relatifs. J'y reviendrai par la suite.

On se pose la question quand le 2-foncteur associé à $f : S \fl T$
\[
f^* :\ \uuGrp(T) \ \fl\ \uuGrp(S)
\]
\printedpage{7}
% head: A. Grothendieck : Dérivateurs, Ch. VIII, Édition des Universités de Montpellier II et Paris VII
% foot: Version 2 février 2002
est 0-fidèle, 1-fidèle, 2-fidèle (\footnote[2]{La $i$-fidélité de $f^*$ ($i \in \{0,1,2\}$) équivaut à celle des foncteurs (obtenus en prenant $G = e_T$)
\[
\underline{\Gamma}(H/T) \ \fl\ \underline{\Gamma}(f^*H/S).
\]}). Si $G$, $H$ sont des fibrés en groupoïdes sur $T$, on considère
\[
\uHom_{\uuGrp(T)}(G,H) \ \fl\ \uHom_{\uuGrp(S)}(f^*G, f^*H),
\]
et suivant que ce foncteur est toujours fidèle, ou pleinement fidèle, ou une équivalence de catégories, on dit que $f^*$ est 0-, 1-, ou 2-fidèle. Dans le cas 2-fidèle, cela implique que si $K$ est un fibré en groupoïdes sur $S$, s'il est équivalent à un $f^*G$, alors $G$ est déterminé à $T$-équivalence près, elle même déterminée à isomorphisme unique près.

\msp{page 8}
Soit $i \in \mathbf{N}$. On dit que $f$ est $i$-asphérique si $\forall t \in T$, $Z = S/t$ est $i$-asphérique, i.e.
\[
\pi_0(Z),\ \dots,\ \pi_i(Z) \qquad\quad \text{triviaux.}
\tag*{$(*)$}
\]

On définit aussi $(-1)$-asphérique : les $S/t$ sont non-vides. De même pour la notion de $i$-coasphérique.

En fait, pour $i \in [-1,+\infty[ \ \subseteq \mathbf{Z}$, on définit un dérivateur $\mathbf{D}_i$ des $i$-types d'homotopie relatifs (cf. plus bas). D'où les notions d'objet $\mathbf{D}_i$-asphérique ($=$ $i$-asphérique), de $\mathbf{D}_i$-équivalence, de foncteur $\mathbf{D}_i$-asphérique et $\mathbf{D}_i$-coasphérique, de foncteur $\mathbf{D}_i$-lisse et $\mathbf{D}_i$-propre, de $\mathbf{D}_i$-fibration et de $\mathbf{D}_i$-fibration \textit{triviale} (i.e. à fibres $\mathbf{D}_i$-asphériques). On doit avoir le

\textbf{Théorème 3.} \textit{Soit $f : S \fl T$ dans} (Cat)\textit{, $i \in \{0,1,2\}$. Supposons l'une des conditions suivantes satisfaite :}

\begin{itemize}[leftmargin=3em]
\item[\textit{a)}] \textit{$f$ est $(i-1)$-propre, et $(i-1)$-asphérique (s'exprime par la condition que certaines catégories, telles les $S/t$, soient $(i-1)$-asphériques).}
\item[\textit{b)}] \textit{$f$ est $(i-1)$-lisse, et $(i-1)$-coasphérique.}
\end{itemize}

\textit{Alors :}

\begin{itemize}[leftmargin=3em]
\item[\textit{1\textsuperscript{o})}] \textit{Le foncteur $f^*$ image inverse}
\[
f^* : \uGrp(T) \ \fl\ \uGrp(S)
\]
\textit{est $i$-fidèle.}
\item[\textit{2\textsuperscript{o})}] \textit{Dans le cas $i = 2$ ($f$ pleinement fidèle), on}
\printedpage{8}
% head: A. Grothendieck : Dérivateurs, Ch. VIII, Édition des Universités de Montpellier II et Paris VII
% foot: Version 2 février 2002
\msp{page 9}
\textit{a de plus ceci : si $G \in \mathrm{Ob}\,\uGrp(S)$ est un fibré en groupoïdes sur $S$, alors $G$ est équivalent à un fibré en groupoïdes de la forme $f^*(H)$ si et seulement si pour tout $t \in T$, la restriction $G_t$ de $G$ à $S_t$ est équivalente à un fibré en groupoïdes constant sur $S_t$. (}\textbf{NB} \textit{La condition a) ou b) implique déjà que $S_t$ est} 1\textit{-connexe, i.e. que son groupoïde fondamental est équivalent à $e$.)}
\item[\textit{3\textsuperscript{o})}] \textit{Soit $Z$ un objet} 1\textit{-connexe dans} (Cat)\textit{, $G$ un fibré en groupoïdes sur $Z$. Pour que $G$ soit $Z$-équivalent à un groupoïde constant sur $Z$, il faut et il suffit que pour toute flèche $u$ dans $Z$, le morphisme de groupoïdes $u^*$ induit sur les fibres de $G$ soit une $Z$-équivalence} (\footnote[3]{On dit alors que $G$ est \textit{essentiellement localement constant} sur $Z$.}).
\end{itemize}

\textbf{Corollaire.} \textit{La condition nécessaire et suffisante de 2\textsuperscript{o}) s'écrit aussi ainsi : pour toute flèche $u$ de $S$ ‘verticale’ sur $T$, i.e. telle que $f(u)$ soit une flèche identique, le foncteur $u^*$ induit par $u$ sur les fibres de $G$ soit une équivalence de groupoïdes.}

% ---- from Der-ChIX-678_p001-005.tex
\printedpage{[non numérotée]}
% head: (aucun titre courant imprimé)

\begin{center}
{\LARGE\textsf{LES DÉRIVATEURS}}

\bigskip

\textbf{ALEXANDRE GROTHENDIECK}

\bigskip\bigskip

\textbf{Chapitre IX}

\textbf{Retour sur les catégories de fractions $\boldsymbol{\mathcal{M}W^{-1}}$,\\
comparaison avec Quillen}
\end{center}

\clearpage

\printedpage{[non numérotée]}
% head: (aucun titre courant imprimé)

Ce texte a été déchiffré et transcrit en \LaTeXe{} par M. Künzer. Il a été édité par M. Künzer, J. Malgoire, et G. Maltsiniotis. La transcription est aussi fidèle que possible au manuscrit. Pour les quelques corrections évidentes, ou rares commentaires des éditeurs, ainsi que pour la numérotation originale des pages du manuscrit, les caractères de machine à écrire \texttt{[typewriter]} entre crochets sont utilisés. Un point d'interrogation entre crochets signifie que l'on n'est pas sûr du mot qui précède.

Cette édition est provisoire. Les remarques, commentaires et corrections sont bienvenus. Envoyer un message à :

\textbf{[e-mail address removed]}

\hfill G. Maltsiniotis\hspace*{2cm}

\clearpage

\printedpage{1}
% head: A. Grothendieck : Dérivateurs, Ch. IX, Édition des Universités de Montpellier II et Paris VII
% foot: Version 12 mars 2002

\msp{page 1}

\section*{1\quad Fonctorialité de $\Ch^{\Sigma}_{\infty}(a,b)$ en $a$, $b$ dans $\mathcal{M}\Sigma^{-1}$}

On se donne une catégorie $\mathcal{M}$ avec \texttt{[un]} ensemble admissible de flèches $\Sigma \subseteq \mathrm{Fl}(\mathcal{M})$, d'où (en prenant p.~ex. $\underline{t} = (t_n)$, avec $t_n = (\underbrace{\Delta_+\Delta_-}_{\tau_0})^{2n}$, et $\underbrace{t_{n+1}}_{=\,\tau_0^{2n+2}} \longrightarrow \underbrace{t_n}_{=\,\tau_0^{2n}}$ défini comme $p * \mathrm{id}_{\tau_0^{2n}} * p$, où $p : \tau_0 \longrightarrow 1$ la dégénérescence.) On a alors une théorie des $\underline{\Ch}_{\Sigma,\infty}(\mathcal{M},a,b)$, que je note simplement $\Ch_\infty(a,b)$ en tant qu'objet de Hot. Des accouplements
\[
\Ch_\infty(a,b) \times \Ch_\infty(b,c) \longrightarrow \Ch_\infty(a,c) \qquad \text{dans Hot}
\]
avec (nous allons l'admettre) \textit{associativité} dans Hot. Cela définit donc des accouplements sur les $\pi_0(\Ch_\infty(a,b)) \overset{\text{déf}}{=} \pi_0(a,b)$, qui n'est autre que la composition des flèches dans $\mathcal{M}\Sigma^{-1}$. D'autre part, on a pour $a,b \in \mathrm{Ob}\,\mathcal{M}$
\[
\Hom_{\mathcal{M}}(a,b) \longrightarrow \pi_0(\Ch_\infty(a,b)) = \pi_0(a,b),
\]
compatible avec les compositions. Mais on peut l'ignorer pour l'instant, et regarder $\mathrm{Ob}\,\mathcal{M}$, ou simplement $\mathcal{M}_0$, avec la seule structure
\[
\left\{
\begin{array}{ccccccc}
\mathcal{M}_0 \times \mathcal{M}_0 & \longrightarrow & \Hot & \qquad & (a,b) & \longmapsto & \Ch_\infty(a,b) \\[1ex]
\mathcal{M}_0 \times \mathcal{M}_0 \times \mathcal{M}_0 & \longrightarrow & \mathrm{Fl}\,\Hot & \qquad & (a,b,c) & \longmapsto & (\Ch_\infty(a,b) \times \Ch_\infty(b,c) \underbrace{\longrightarrow}_{\substack{\text{flèche}\\ \text{dans Hot}}} \Ch_\infty(a,c))
\end{array}
\right.
\]
avec condition \texttt{[d']}associativité, et aussi une condition d'unités :
\[
\Ch_\infty(a,a) \text{ est muni d'un } 1_a \in \pi_0\Ch_\infty(a,a) = \pi_0(a,a),
\]
jouant le rôle d'une `identité bilatère', dans le sens suivant. Notons que les accouplements
\[
\begin{array}{ccccc}
\Ch_\infty(a,b) & \times & \Ch_\infty(b,c) & \longrightarrow & \Ch_\infty(a,c) \\[1ex]
P & \times & Q & \longrightarrow & R
\end{array}
\]
\msp{page 2}
définissent
\[
\begin{array}{ccc}
P & \longrightarrow & \underline{\Hom}_{\Hot}(Q,R) \\[1ex]
Q & \longrightarrow & \underline{\Hom}_{\Hot}(P,R),
\end{array}
\]
d'où en passant aux $\pi_0$
\[
\begin{array}{ccc}
\pi_0(P) & \longrightarrow & \Hom_{\Hot}(Q,R) \\[1ex]
\pi_0(Q) & \longrightarrow & \Hom_{\Hot}(P,R),
\end{array}
\]
\printedpage{2}
% head: A. Grothendieck : Dérivateurs, Ch. IX, Édition des Universités de Montpellier II et Paris VII
% foot: Version 12 mars 2002
en d'autres termes,
\[
\begin{array}{ccc}
\pi_0(a,b) & \longrightarrow & \Hom_{\Hot}(\Ch_\infty(b,c),\Ch_\infty(a,c)) \\[1ex]
\pi_0(b,c) & \longrightarrow & \Hom_{\Hot}(\Ch_\infty(a,b),\Ch_\infty(a,c)).
\end{array}
\]

On en déduit en particulier pour $a,b \in \mathcal{M}_0$
\[
\left\{
\begin{array}{cccc}
1) & \pi_0(a,a) & \longrightarrow & \Hom_{\Hot}(\Ch_\infty(a,b),\Ch_\infty(a,b)) \\[1ex]
2) & \pi_0(a,a) & \longrightarrow & \Hom_{\Hot}(\Ch_\infty(b,a),\Ch_\infty(b,a)),
\end{array}
\right.
\tag*{$(*)$}
\]
et un objet $\alpha_a$ de $\pi_0(a,a)$ est unité bilatère si ses deux images sont les homomorphismes identiques dans $\Ch_\infty(a,b)$ resp. $\Ch_\infty(b,a)$. Passant aux $\pi_0$ de ceux-ci, cela implique que $\alpha_a$ est unité bilatère pour la composition des $\pi_0(x,y)$, relativement à $b$. Si on veut que ce soit le cas pour \textit{tout $b$,} cela signifie (prenant $b = a$) que $\alpha_a$ est un objet unité bilatère pour la structure de composition associative (donc de monoïde) dans $\pi_0(a,a)$, donc que $\pi_0(a,a)$ est un monoïde \textit{unitaire,} ayant $\alpha_a$ pour objet unité. L'élément $\alpha_a$ est donc déterminé de façon unique. Et on veut que $\alpha_a$ existe
\msp{page 3}
pour tout $a \in \mathcal{M}_0$, donc $\mathcal{M}_0$, muni des $\pi_0(a,b)$ et de leur composition, devient une \textit{catégorie,} soit $\mathcal{M}_\Pi$; et les flèches identiques de celle-ci doivent jouer le rôle d'unité bilatère également pour les `Hom externes' $\Ch_\infty(a,b)$, à valeurs dans Hot.

On peut dire que pour $a = a_0$ fixé, les $\Ch_\infty(a_0,b)$ ($b$ variable) forment les valeurs d'un \textit{foncteur} covariant
\[
\mathcal{M}_\Pi \longrightarrow \Hot \qquad b \longmapsto \Ch_\infty(a_0,b),
\]
en associant à
\[
b \xrightarrow{\ u\ } b', \qquad u \in \pi_0(b,b'),
\]
la flèche
\[
\Ch_\infty(a,b) \longrightarrow \Ch_\infty(b') % [sic] \Ch_\infty(b') -> \Ch_\infty(a_0,b'), argument manquant (cf. la note des éditeurs qui suit)
\]
\texttt{[plutôt} $\Ch_\infty(a_0,b) \longrightarrow \Ch_\infty(a_0,b')$\texttt{]} correspondante $*_1$ [?]. Que ce soit bien là une loi fonctorielle doit résulter formellement de l'associativité des accouplements sur les $\Ch_\infty$. De même, on a une loi fonctorielle, pour $b$ fixé,
\[
\mathcal{M}^o_\Pi \longrightarrow \Hot, \qquad a \longmapsto \Ch_\infty(a,b),
\]
grâce à l'associativité. En fait, l'associativité impliquera aussi qu'on a un \textit{bifoncteur}
\[
\begin{array}{ccc}
\mathcal{M}^o_\Pi \times \mathcal{M}_\Pi & \longrightarrow & \Hot \\[1ex]
(a\,,\,b) & \longmapsto & \Ch_\infty(a,b),
\end{array}
\]
\printedpage{3}
% head: A. Grothendieck : Dérivateurs, Ch. IX, Édition des Universités de Montpellier II et Paris VII
% foot: Version 12 mars 2002
\msp{page 4}
car si
\[
a' \xrightarrow{\ u\ } a, \qquad b \xrightarrow{\ v\ } b',
\]
on a le diagramme
\[
\begin{tikzcd}[column sep=large, row sep=large]
\Ch_\infty(a,b) \arrow[r, "\substack{\text{comp. avec}\\ v}"] \arrow[d, "\text{comp. avec } u"'] & \Ch_\infty(a,b') \arrow[d, "\text{comp. avec } u"] \\
\Ch_\infty(a',b) \arrow[r, "\substack{\text{comp. avec}\\ v}"] & \Ch_\infty(a',b'),
\end{tikzcd}
\]
qui est commutatif grâce à l'associativité pour
\[
\Ch_\infty(a',a),\ \Ch_\infty(a,b),\ \Ch_\infty(b,b').
\]

\textbf{Corollaire} de cette dépendance bifonctorielle : \textit{Si $u : a' \longrightarrow a$, $v : b \longrightarrow b'$ sont dans $\mathcal{M}_\Pi$, alors la flèche correspondante}
\[
\Ch_\infty(a,b) \longrightarrow \Ch_\infty(a',b') \qquad \textit{dans } \Hot
\]
\textit{est un} isomorphisme.

Les mêmes arguments devraient marcher si au lieu des $\Hom\,\mathrm{ext}$ à valeurs dans Hot, on les prenait dans Ind\,Hot, ou dans Cat\,hot, ou dans Ind\,Cat\,hot -- je pense bien sûr aux $\mathrm{CH}_\infty(a,b)$ (dans la théorie de la localisation) dans Ind\,Cat\,hot, qui peuvent être regardés comme provenant d'un bifoncteur (où $\mathcal{M}_\Pi = \mathcal{M}\Sigma^{-1}$)
\[
\mathcal{M}_\Pi \times \mathcal{M}_\Pi \longrightarrow \mathrm{Ind}\,\mathrm{Cat}\,\mathrm{hot}.
\]

\msp{page 5}
Voici l'axiomatisation définitive (?). On remplace Hot par une catégorie

\begin{itemize}[label={$\mathcal{H}$,}]
\item catégorie avec produit $*$ associatif, commutatif (?) (\footnote{On ne doit pas avoir besoin de la commutativité.}), unitaire, et toutes les compatibilités habituelles au besoin,
\end{itemize}

l'objet unité pour $*$ sera noté 1, et on notera
\[
\left\{
\begin{array}{ccc}
\pi_0(\xi) & \overset{\text{déf}}{=} & \Hom_{\mathcal{H}}(1,\xi) \\[1ex]
\pi_0 : \mathcal{H} & \longrightarrow & \mathrm{Ens}.
\end{array}
\right.
\]

Si on a
\[
P * Q \xrightarrow{\ u\ } R,
\]

% ---- from Der-ChIX-678_p006-010.tex
\printedpage{4}
% head: A. Grothendieck : Dérivateurs, Ch. IX, Édition des Universités de Montpellier II et Paris VII
% foot: Version 12 mars 2002

on en déduit
\[
\left\{
\begin{array}{lcl}
\pi_0(P) & \xrightarrow{c_u^{\mathrm{l}}} & \mathrm{Hom}_{\mathcal{H}}(Q,R)\\
\pi_0(Q) & \xrightarrow{c_u^{\mathrm{r}}} & \mathrm{Hom}_{\mathcal{H}}(P,R)
\end{array}
\right.
\]
ainsi :
\[
\begin{array}{rlcl}
\text{Si } \alpha\in\pi_0(P),\ \text{i.e. } 1\xrightarrow{\ \alpha\ }P,\ \text{alors} & \multicolumn{3}{l}{\alpha*\mathrm{id}_Q : Q\longrightarrow P*Q}\\
\text{et} & c^{\mathrm{l}}_u(\alpha) & = & u\circ(\alpha*\mathrm{id}_Q)\\
\text{et dualement} & c^{\mathrm{r}}_u(\beta) & = & u\circ(\mathrm{id}_P*\beta).
\end{array}
\]
Si on a trois objets $P_0$, $P_1$, $P_2$ dans $\mathcal{H}$, des accouplements
\[
\left\{
\begin{array}{rcl}
P_0*P_1 & \xrightarrow{u_1} & Q_0\\
P_1*P_2 & \xrightarrow{u_2} & Q_1\\
P_0*Q_1 & \xrightarrow{u_3} & R\\
Q_0*P_2 & \xrightarrow{u_4} & R
\end{array}
\right.
\]
\[
\begin{tikzcd}
a \arrow[r, "P_0"'] \arrow[rr, bend left=20, "Q_0"] \arrow[rrr, bend left=35, "R"] & b \arrow[r, "P_1"'] \arrow[rr, bend right=20, "Q_1"'] & c \arrow[r, "P_2"'] & d
\end{tikzcd}
\]

\msp{page 6}
avec une relation d'associativité, disant que
\[
\begin{tikzcd}
P_0*P_1*P_2 \arrow[r] \arrow[d] & Q_0*P_2 \arrow[d]\\
P_0*Q_1 \arrow[r] & R
\end{tikzcd}
\]
est commutatif, alors on en déduit deux relations de transitivité et une de commutation, pour les opérations associées aux éléments de $\pi_0(P_0)$, $\pi_0(P_1)$, $\pi_0(P_2)$, qu'on n'explicitera pas, sauf l'une :
\[
\begin{array}{rrrcl}
\multicolumn{5}{l}{\text{Soit } \beta\in P_1,\ \gamma\in P_2,}\\
\text{d'où par } \beta : & c_{u_1}(\beta) : & P_0 & \longrightarrow & Q_0,\\
\text{et par } \gamma : & c_{u_4}(\gamma) : & Q_0 & \longrightarrow & R,\\
\text{et par } \beta *_{u_2}\gamma\in\pi_0(Q_1) : & c_{u_3}(\beta\,\texttt{[}{*_{u_2}}\texttt{]}\gamma) : & P_0 & \longrightarrow & R,
\end{array}
\]

\printedpage{5}
% head: A. Grothendieck : Dérivateurs, Ch. IX, Édition des Universités de Montpellier II et Paris VII
% foot: Version 12 mars 2002
\texttt{[plutôt} $\beta\in\pi_0(P_1)$\texttt{,} $\gamma\in\pi_0(P_2)$\texttt{]}, et bien on a
\[
\boxed{\;c(\beta *_{u_2}\gamma) \;=\; c(\gamma)c(\beta).\;}
\]
Il faut encore définir $\beta *_{u_2}\gamma$, via l'accouplement $P_1*P_2\xrightarrow{u_2}Q_1$, on a en effet
\[
\begin{array}{ccc}
\pi_0(P_1)\times\pi_0(P_2) & \longrightarrow & \pi_0(P_1*P_2)\\
(\beta\,,\,\gamma) & \longmapsto & \beta*\gamma,
\end{array}
\]
et on compose avec
\[
\pi_0(u_2) : P_1*P_2 \longrightarrow Q_1
\]
\texttt{[plutôt} $\pi_0(u_2) : \pi_0(P_1*P_2) \longrightarrow \pi_0(Q_1)$\texttt{]}.

\msp{page 7}
On se donne un ensemble $\mathcal{M}_0$, et
\[
\left\{
\begin{array}{rclccrcl}
\mathcal{M}_0\times\mathcal{M}_0 & \xrightarrow{\mathrm{Ch}} & \mathcal{H} & \text{noté} & & (a,b) & \longmapsto & \mathrm{Ch}(a,b)\ \text{ou}\ \mathrm{Hom\,ext}_{\mathcal{H}}(a,b)\\
\mathcal{M}_0\times\mathcal{M}_0\times\mathcal{M}_0 & \xrightarrow{\text{comp.}} & \mathrm{Fl}\,\mathcal{H} & & & (a,b,c) & \longmapsto & \bigl(\mathrm{Ch}(a,b)\times\mathrm{Ch}(b,c)\longrightarrow\mathrm{Ch}(a,c)\bigr)
\end{array}
\right.
\]
avec deux axiomes

\begin{enumerate}
\item[a)] Associativité pour les compositions, pour $(a,b,c,d)\in\mathcal{M}_0^4$
\[
\begin{tikzcd}
\mathrm{Ch}(a,b)\times\mathrm{Ch}(b,c)\times\mathrm{Ch}(c,d) \arrow[r] \arrow[d] & \mathrm{Ch}(a,c)\times\mathrm{Ch}(c,d) \arrow[d]\\
\mathrm{Ch}(a,b)\times\mathrm{Ch}(b,d) \arrow[r] & \mathrm{Ch}(a,d).
\end{tikzcd}
\]
\item[b)] Existence d'objets unité bilatères
\[
1_a\in\pi_0(\mathrm{Ch}(a,a)) \qquad \text{pour tout } a\in\mathcal{M}_0
\]
\end{enumerate}

Alors les $\pi_0(\mathrm{Ch}(a,b))$ et leur composition définissent sur $\mathcal{M}_0$ comme ensemble d'objets une structure de \textit{catégorie} $\MOP$, ayant les $1_a$ comme identités, et on trouve que
\[
(a,b)\longmapsto\mathrm{Ch}(a,b)
\]
peut s'interpréter à partir d'un \textit{bifoncteur} canonique
\[
\mathcal{M}^o_{0\,\Pi}\times\MOP \longrightarrow \mathcal{H},
\]
dont le composé avec
\[
\mathcal{H}\xrightarrow{\pi_0}\mathrm{Ens}
\]
n'est autre que le bifoncteur $(a,b)\longrightarrow\mathrm{Hom}_{\MOP}(a,b)$,
\printedpage{6}
% head: A. Grothendieck : Dérivateurs, Ch. IX, Édition des Universités de Montpellier II et Paris VII
% foot: Version 12 mars 2002
\msp{page 8}
\[
\mathrm{Hom}_{\MOP}(a,b) \;\simeq\; \pi_0(\mathrm{Ch}(a,b)).
\]
Cette relation est la \textit{définition} de $\underline{\mathrm{Hom}}_{\MOP}$, mais ici la relation est entendue comme une % ?? uncertain: underline under Hom -- single underline or double underline
bijection \textit{fonctorielle} en $a$, $b$.

On trouve notamment que si
\[
u : a'\longrightarrow a, \qquad v : b\longrightarrow b'
\]
sont des flèches \textit{inversibles} dans $\MOP$, alors
\[
\mathrm{Ch}(a,b)\longrightarrow\mathrm{Ch}(a',b')
\]
correspondant dans $\mathcal{H}$ est \texttt{[un]} \textit{isomorphisme.}

On appliquera tout ceci surtout aux deux situations déduites d'une situation $(\mathcal{M},\Sigma)$, avec
\[
\begin{array}{rlcllcl}
1^o) & \mathcal{H} & = & \mathrm{Hot}, & \mathrm{Ch}(a,b) & = & \mathrm{Ch}^{\Sigma}_{\infty}(\mathcal{M};a,b)\\
2^o) & \mathcal{H} & = & \mathrm{Ind\,Cat\,hot}, & \mathrm{Ch}(a,b) & = & \mathrm{CH}^{\Sigma}_{\infty}(\mathcal{M};a,b)\ \text{(\footnotemark[2])}.
\end{array}
\]
\footnotetext[2]{$*$ dans Hot et Ind\,Cat\,hot est le produit cartésien sans plus.}
Dans les deux cas, on a donc un foncteur sur $(\mathcal{M}\Sigma^{-1})^o\times(\mathcal{M}\Sigma^{-1})$, à valeurs soit dans Hot, soit dans Ind\,Cat\,hot. (\textbf{NB} Les deux structures de catégorie sur $\mathcal{M}_0$ déduites des % [sic] des -> de, « déduites des ces deux foncteurs »
ces deux foncteurs sont canoniquement isomorphes.)
\[
\begin{tikzcd}[column sep=small]
 & \mathrm{Hot} & (a,b)\longmapsto\mathrm{Ch}^{\Sigma}_{\infty}(a,b)\\
(\mathcal{M}\Sigma^{-1})^o\times(\mathcal{M}\Sigma^{-1}) \arrow[ur] \arrow[dr] & & \\
 & \mathrm{Ind\,Cat\,hot} & (a,b)\longmapsto\mathrm{CH}^{\Sigma}_{\infty}(a,b)
\end{tikzcd}
\]

\msp{page 9}
Comme les foncteurs canoniques
\[
\begin{tikzcd}
\mathrm{Ind\,Cat\,hot} \arrow[r, "\mathrm{Ind}\,\varphi"] & \mathrm{Ind\,Hot} \arrow[r, "\mathrm{Ind}\,\psi"] & \mathrm{Ind\,Hot}_0^{\wedge} \arrow[r, "\varinjlim"] & \mathrm{Hot}_0^{\wedge},\\
\mathrm{Cat\,hot} \arrow[u, hook] \arrow[r, "\varphi"] & \mathrm{Hot} \arrow[u, hook] \arrow[r, "\psi"] & \mathrm{Hot}_0^{\wedge} \arrow[u, hook] \arrow[ur, "\mathrm{id}"'] &
\end{tikzcd}
\]
commutant aux produits finis, toute théorie de `Hom externes' à valeurs dans une de ces catégories, en définit une à valeurs dans chacune de celles dans laquelle elle s'envoie.
\printedpage{7}
% head: A. Grothendieck : Dérivateurs, Ch. IX, Édition des Universités de Montpellier II et Paris VII
% foot: Version 12 mars 2002
D'ailleurs, ces foncteurs ne sont pas tous pleinement fidèles (c'est en tous cas faux pour $\varphi$ et $\mathrm{Ind}\,\varphi$, et j'ignore ce qui en est de $\psi$ et $\mathrm{Ind}\,\psi$ -- les flèches verticales par contre sont pleinement fidèles (\footnote[3]{et également, $\mathrm{Ind\,Hot}_0^{\wedge}\xrightarrow{\varinjlim}\mathrm{Hot}_0^{\wedge}$ n'est \textit{pas} pleinement fidèle.}).) Mais ils induisent des bijections pour les $\pi_0$. Donc la catégorie $\MOP$ déduite d'une théorie de Hom externes à valeurs dans une de ces catégoires, et celle % [sic] catégoires -> catégories, faute d'orthographe
déduite dans une autre où elle s'envoie, est la même à isomorphisme canonique près. Pour ce qui est de cette définition de $\MOP$, rien n'est changé donc quand on regarde les $\mathrm{CH}_{\infty}$ comme des objets de Ind\,Hot, ou seulement de $\mathrm{Ind\,Hot}_0^{\wedge}$, ou même seulement de $\HotO$, ou quand on regarde les $\mathrm{Ch}_{\infty}$ comme des objets
\msp{page 10}
dans $\HotO$. Mais bien sûr, la donnée des accouplements de composition perd chaque fois en précision (sauf peut-être en passant de Hot à $\HotO$, si $\psi$ est pleinement fidèle \dots\,). Quand on prend les images des $\mathrm{CH}^{\Sigma}_{\infty}(a,b)$ et $\mathrm{Ch}^{\Sigma}_{\infty}(a,b)$ dans $\HotO$ (la dernière catégorie du diagramme), on trouve des objets canoniquement isomorphes (fonctoriellement en $a$, $b$ dans $\MOP\simeq\mathcal{M}\Sigma^{-1}$). Quand on les regarde comme étant respectivement dans Ind\,Hot et dans $\mathrm{Hot}\hookrightarrow\mathrm{Ind\,Hot}$, on trouve un homomorphisme fonctoriel
\[
\mathrm{CH}^{\Sigma}_{\infty}(a,b)\longrightarrow\mathrm{Ch}^{\Sigma}_{\infty}(a,b) \qquad\quad \text{dans Ind\,Hot ou dans } \mathrm{Ind\,Hot}_0^{\wedge} \text{ au choix.} \tag*{$(*)$}
\]
La $\varinjlim$ de $\mathrm{CH}^{\Sigma}_{\infty}(a,b)$ dans Hot n'existe pas en général (sauf si l'ind-objet est essentiellement constant), mais elle existe dans $\HotO$, et l'isomorphisme canonique
\[
\varinjlim_{\HotO}\mathrm{CH}^{\Sigma}_{\infty}(a,b)\xrightarrow{\ \sim\ }\mathrm{Ch}^{\Sigma}_{\infty}(a,b) \qquad\quad \text{dans } \HotO \tag*{$(**)$}
\]
est la flèche déduite de $(*)$ dans $\mathrm{Ind\,Hot}_0^{\wedge}$.

\msp{page 11}
Dans les deux cas CH et Ch (qu'on va désigner par le symbole commun $C_{\infty}(a,b)$), la fonctorialité de cet objet $C_{\infty}(a,b)$ (même au sens le plus faible, comme objet de $\HotO$) par rapport à $a$ et $b$ (même au sens le plus faible, comme foncteur $\mathcal{M}^o\times\mathcal{M}\longrightarrow\HotO$ et pas même $(\mathcal{M}\Sigma^{-1})^o\times(\mathcal{M}\Sigma^{-1})\longrightarrow\HotO$) n'est pas évidente du tout sur la définition directe en termes du système inductif explicite dans Cat
\[
C_{\infty}(a,b) \;=\; (\mathrm{Ch}_{\tau_0^{2n}}(a,b))_{n\in\mathbf{N}},
\]
car une flèche $a'\longrightarrow a$ ou $b\longrightarrow b'$ dans $\mathcal{M}$, et même qu'elle serait dans $\Sigma$, ne définit aucunement une flèche dans Cat
\[
\begin{array}{rcl}
C_{\infty}(a,b) & \longrightarrow & C_{\infty}(a',b)\\
\text{ou}\ \ C_{\infty}(a,b) & \longrightarrow & C_{\infty}(a,b').
\end{array}
\]
Cette difficulté rappelle celle qu'avait à surmonter \textsc{Quillen}, dans Homotopical Algebra, pour affronter la différence entre `left homotopies' et `right homotopies'. Je me rappelle comment
\printedpage{8}
% head: A. Grothendieck : Dérivateurs, Ch. IX, Édition des Universités de Montpellier II et Paris VII
% foot: Version 12 mars 2002
\msp{page 12}
on peut néanmoins définir de telles flèches, comme résultant d'une flèche (beaucoup plus riche) de `composition' de chemins entre $a'$ et $a$, et $a$ et $b$, ou entre $a$ et $b$, et $b$ et $b'$, soit en général donc l'accouplement
\[
C_{\infty}(a,b)\times C_{\infty}(b,c)\longrightarrow C_{\infty}(a,c),
\]
dont la définition directe n'est pas moins problématique, et ne peut en tous cas \textit{pas} se donner comme une flèche dans Cat lui-même, ou dans Ind\,Cat. Interprétant p.~ex. $\mathrm{Ch}_{\infty}(a,b)$ en termes de chemins \textit{infinis} (des deux cotés) de $a$ à $b$, en un sens évident (mais % [sic] cotés -> côtés, accent omis
`essentiellement finis' des deux cotés), on peut définir aussi % [sic] cotés -> côtés, accent omis
\[
\mathrm{Ch}^{-}_{\infty}(a,b), \qquad \mathrm{Ch}^{+}_{\infty}(a,b)
\]
comme les sous-catégories pleines formées respectivement de chemins qui sont totalement dégénérés en degrés $\geq 0$ (i.e. $a_0\longleftarrow a_1\longrightarrow a_2\longleftarrow a_3\cdots$ sont tous des identités), ce qui revient aux chemins infinis seulement à gauche, i.e. s'arrêtant à $a_0$, et symétriquement pour $\mathrm{Ch}^{+}_{\infty}(a,b)$. Même tabac pour les $\mathrm{CH}^{+}$, $\mathrm{CH}^{-}$.

\msp{page 13}
Ceci dit, on a ceci (revenant à la notion $C_{\infty}$ ambivalente) :

\begin{enumerate}
\item[a)] les inclusions pleinement fidèles
\[
C^{-}_{\infty}(a,b)\hookrightarrow C_{\infty}(a,b)\hookleftarrow C^{+}_{\infty}(a,b)
\]
sont des isomorphismes dans Ind\,Cat\,hot resp. dans Hot, et
\item[b)] On a bel et bien une composition des chemins
\[
C^{-}_{\infty}(a,b)\times C^{+}_{\infty}(b,c)\longrightarrow C_{\infty}(a,c).
\]
\end{enumerate}

L'\textit{associativité} dans le cas du CH n'offre d'ailleurs pas de difficulté, à cause de la possibilité de se ramener aux
\[
C_m(a,b)\times C_n(b,c)\times C_p(c,d)\rightrightarrows C_N(a,d)
\]
pour $N$ grand, et prouver l'egalité comme flèches dans Cat\,hot -- ce qui résultera du lemme % [sic] egalité -> égalité, accent omis
d'homotopie sur les intervalles dans Cat. Dans le cas du $\mathrm{Ch}_{\infty}$ par contre je n'ai pas trouvé de démonstration. (Faute seulement d'avoir essayé sérieusement de m'y mettre? Je n'ai pas trop envie, vu qu'il se pourrait que $\mathrm{Hot}\longrightarrow\HotO$ soit pleinement fidèle, auquel cas le travail devient inutile \dots\,)

% ---- from Der-ChIX-678_p011-015.tex
\printedpage{9}
% head: A. Grothendieck : Dérivateurs, Ch. IX, Édition des Universités de Montpellier II et Paris VII
% foot: Version 12 mars 2002

\msp{page 14}

\section*{2\quad Comparaison avec Quillen}

Il se pose un certain nombre de questions sur ce foncteur $(\M\Sigma^{-1})^o \times (\M\Sigma^{-1}) \longrightarrow \Hot$, et sa relation aux constructions classiques, notamment à la théorie de \textsc{Quillen}. Je vais maintenant noter
\[
\begin{aligned}
\Ch^{\Sigma}_{\infty}(a,b) \;&=\; \Homext_{\Hot}(a,b) && \text{ou simplement}\\
&\phantom{\;=\;} \Homext(a,b).
\end{aligned}
\]

1$^o$). Soient $f,g : a \rightrightarrows b$ dans $\M$, alors (si $\M$ est une `model category') \textsc{Quilllen} définit un ensemble % [sic] Quilllen -> Quillen, three l's printed
\[
\pi_1(a,b;f,g) \qquad \text{ou} \qquad \pi_1(f,g)
\]
des classes d'homotopie de $f$ à $g$, de telle façon que $\mathrm{Hom}_{\M}(a,b)$ devient l'ensemble des objets d'un groupoïde, dont les flèches sont les éléments de $\pi_1(f,g)$, pour une notion de composition naturelle. Je suspecte que l'on a
\[
\pi_1(f,g) \;=\; \mathrm{Hom}_{\underbrace{\Pi(\uCh^{\Sigma}_{\infty}(a,b))}_{\substack{\text{groupoïde}\\ \text{fodamental}\\ \text{de } \uCh^{\Sigma}_{\infty}(a,b)}}}(f,g), % [sic] fodamental -> fondamental, missing n
\]
quand $f$ et $g$ sont considérés comme des `chemins infinis' (en insérant chaque $f : a \longrightarrow b$ dans
\[
\cdots \underbrace{a_{-2}}_{=a} \xleftarrow{\ \mathrm{id}\ } \underbrace{a_{-1}}_{=a} \xrightarrow{\ f\ } \underbrace{a_{0}}_{=b} \xleftarrow{\ \mathrm{id}\ } \underbrace{a_{1}}_{=b} \longrightarrow \cdots\ \ )
\]

\msp{page 15}

Donc la définition des $\pi_1(f,g)$ (si c'est bien ça) se prolonge de façon naturelle aux chemins infinis (essentiellement finis) arbitraires.

2$^o$) Supposons que $\M$ admette un objet à la fois initial et final $0_{\M}$ ($\M$ `ponctuée'). Je dis qu'alors toutes les catégories $\uCh^{\Sigma}_{\infty}(a,b)$ sont ponctuées de façon naturelle : à savoir, par le chemin infini de $a$ à $b$ défini (comme plus haut) à partir de $f = 0_{a,b} : a \longrightarrow b$, l'application nulle de $a$ à $b$. On a envie alors de regarder les $\uCh^{\Sigma}_{\infty}(a,b)$ comme définissant des éléments dans le $\Hot_{\bullet}$ (`Hot ponctué'). C'est plus fin que de dire qu'on se donne une classe neutre
\[
e \;\longrightarrow\; \Ch(a,b) \qquad \text{dans Hot,}
\]
car le foncteur de $\Hot_{\bullet}$ dans Hot (ou dans $e\backslash\Hot$) n'est pas fidèle. Le problème, c'est de définir les flèches
\[
\Ch(a,b)\times\Ch(b,c) \;\longrightarrow\; \Ch(a,c)
\]
dans $\Hot_{\bullet}$, et pas seulement dans Hot, et avec associativité (et les unités) dans $\Hot_{\bullet}$ également. Si on applique la définition avec $\Ch^{-}(a,b)$ et $\Ch^{+}(b,c)$, le composé des éléments nuls $a \xrightarrow{\ 0\ } b$ de $\Ch^{-}(a,b)$ et $b \xleftarrow{\ \mathrm{id}\ } b \xrightarrow{\ 0\ } c$ de $\Ch^{+}(b,c)$
\printedpage{10}
% head: A. Grothendieck : Dérivateurs, Ch. IX, Édition des Universités de Montpellier II et Paris VII
% foot: Version 12 mars 2002
\msp{page 16}
est l'élément de $\Ch_{\infty}(a,c)$ que voici
\[
\begin{tikzcd}[row sep=huge]
 & a_{-2} & a_{-1} \arrow[l] \arrow[r] & a_{0} \arrow[d, equal] & a_{1} \arrow[l] \arrow[r] & a_{2} & \cdots \\
{}_{\mathrm{id}}\ \cdots & a & a \arrow[l, "\mathrm{id}"'] \arrow[r, "0"'] & b & b \arrow[l, "\mathrm{id}"'] \arrow[r, "0"'] & c & \cdots_{\mathrm{id}} \arrow[l, "\mathrm{id}"']
\end{tikzcd}
\]
ce n'est \textit{pas} l'élément nul de $\Ch_{\infty}(a,c)$ ! Mais on a des flèches canoniques
\[
e^{-}_{a,b}\circ e^{+}_{b,c} \;\;\rightleftarrows\;\; e_{a,c} \qquad \text{dans } \Ch_{\infty}(a,c)
\]
qui sont l'identité sur tous les sommets sauf celui d'indice 0. Il y a un sorite à développer, comme quoi une paire
\[
(f : K_{\bullet} \longrightarrow L_{\bullet}, c)
\]
d'un foncteur entre catégories ponctuées, et d'un chemin $c : f(e_{K_{\bullet}}) \longrightarrow e_{L_{\bullet}}$, définit dans Hot une flèche canonique
\[
\mathrm{hot}_{\bullet}(K_{\bullet}) \;\longrightarrow\; \mathrm{hot}_{\bullet}(L_{\bullet})
\]
avec transitivité si on a
\[
(f' : L_{\bullet} \longrightarrow M_{\bullet}, c' : f'(e_{L_{\bullet}}) \longrightarrow e_{M_{\bullet}}),
\]

\msp{page 17}

d'où
\[
f'' = f'\circ f : K_{\bullet} \longrightarrow M_{\bullet}, \qquad \text{et}\ \ c'' : f''(e_{K_{\bullet}}) \longrightarrow e_{M_{\bullet}}), % [sic] closing parenthesis without opening parenthesis -> (f'' = f'\circ f : ..., c'' : ...), unbalanced parenthesis
\]
définit ainsi
\[
\begin{tikzcd}
f'(f(e_{K_{\bullet}})) \arrow[r, "f'(c)"] \arrow[rr, bend right=12, "{c''\ \text{par définition}}"'] & f'(e_{L_{\bullet}}) \arrow[r, "c'"] & e_{M_{\bullet}}
\end{tikzcd}
\]

\noindent\llap{?\quad}Cela nous donne donc bel et bien des accouplements dans $\Hot_{\bullet}$ (et aussi dans $\mathrm{Cat\,hot}_{\bullet}$?). Il est difficile de croire qu'ils pourraient ne pas être associatifs! D'ailleurs $\Hot_{\bullet}$, avec le produit cartésien ordinaire (\footnote[4]{Non, il faut prendre le produit $\xi_{\bullet} \divideontimes \eta_{\bullet}$ ($\xi_{\bullet}\times\eta_{\bullet}$, avec $\xi_{\bullet}\vee\eta_{\bullet}$ contracté en $e$).}), satisfait à l'axiomatique de tantôt pour $\mathcal{H}$, et $\Hot_{\bullet} \longrightarrow \Hot$ est compatible avec ce produit (\footnote[5]{Non, mais on s'en tire quand même \ldots}), et avec $\pi_0$. % ?? uncertain: product symbol in footnote 4 -- \divideontimes (⋇) or reference mark ※ (× with four dots)

Si tout ceci marche, on aurait donc une théorie des $\Homext_{\bullet}$ à valeurs dans $\Hot_{\bullet}$; et sans doute aussi pour les $\mathrm{CH}_{\infty}$ comme objets de $\mathrm{Ind\,Hot}_{\bullet}$, et non seulement de $\mathrm{Ind\,Hot}$. Mais
\printedpage{11}
% head: A. Grothendieck : Dérivateurs, Ch. IX, Édition des Universités de Montpellier II et Paris VII
% foot: Version 12 mars 2002
de toutes façons on ne peut en général laisser tomber Cat hot, qui reste importante -- donc il faudrait avoir des objets $\mathrm{CH}_{\infty}(a,b)$ dans $\mathrm{Ind\,Cat\,hot}_{\bullet}$, pour avoir une côte bien taillée.

Les $\pi_0(C(a,b))$ seront à présent des ensembles ponctués, et ce sera par l'image de l'application nulle $a \xrightarrow{\ 0\ } b$, i.e. le composé
\[
a \longrightarrow 0 \longrightarrow b
\]

\msp{page 18}

qu'il faut composer avec l'elément déduit du chemin tronqué % [sic] l'elément -> l'élément, missing accent
\[
a \longrightarrow \underbrace{0}_{=a_0} \longleftarrow b \qquad \ldots
\]

Il faudrait montrer que 0 est aussi un objet neutre de $\M\Sigma^{-1}$, i.e. que
\[
\pi_0(a,0) = \pi_0(0,a) = \{0\} \qquad \text{pour tout } a \in \mathrm{Ob}\,\M
\]
et interpréter l'élément nul dans $\pi_0(a,b)$ comme l'application nulle dans $\M\Sigma^{-1}$. De façon générale, il faut des conditions qui assurent (sans que $\M$ soit forcément ponctuée) que
\[
\M \xrightarrow{\ \gamma\ } \M\Sigma^{-1}
\]
commute aux sommes finies, aux produits finis (donc $\gamma(\varphi_{\M}) = \varphi_{\M\Sigma^{-1}}$, resp. $\gamma(e_{\M}) = e_{\M\Sigma^{-1}}$) (\footnote[6]{à faire.}).

Supposons enfin que $\M$ soit \textit{additive,} et même $k$-linéaire ($k$ un anneau fixé), i.e. on a de plus un homomorphisme d'anneaux
\[
k \;\longrightarrow\; \mathrm{End}(\mathrm{id}_{\M}).
\]

On doit supposer quelque compatibilité de $\Sigma$ à la structure additive, p. ex. que la somme directe de deux flèches dans $\Sigma$ est dans $\Sigma$, et sans doute quelque chose de plus fort, du genre `lemme des cinq' pour $\Sigma$. (Mais on ne suppose pas forcément $\M$ abélienne, donc on aurait dû mal à formuler un lemme des cinq?) % [sic] dû mal -> du mal, accent on du

\msp{page 19}

Mais peut-être ici faut-il un axiome plus subtil, p.ex. (si $\underbrace{\Sigma}_{\text{suspension}}$ et $\Omega$ existent (\footnote[7]{cf. plus bas ce que ça signifie.})) que l'ensemble des morphismes d'adjonction $\Sigma\Omega \longrightarrow \mathrm{id}$, $\mathrm{id} \longrightarrow \Omega\Sigma$ soit \texttt{[formé d']}isomorphismes, ce qui implique que $\Omega$ resp. $\Sigma$ est pleinement fidèle. En tous cas, l'attente naturelle ici, c'est que $\Ch^{\Sigma}_{\infty}(a,b)$ soit non seulement dans $\Hot_{\bullet}$, mais même dans la catégorie $D^{\leq 0}(k\text{-Mod})$, sous-catégorie pleine de la catégorie dérivée $D^{-}(k\text{-Mod})$ formée des $K$ tels que $H^{i}(K) = 0$ si $i > 0$. On peut aussi la noter $D_{\geq 0}(k\text{-Mod})$, l'interprétant en
\printedpage{12}
% head: A. Grothendieck : Dérivateurs, Ch. IX, Édition des Universités de Montpellier II et Paris VII
% foot: Version 12 mars 2002
termes de complexes de `chaînes' plutôt que de cochaînes. Ici les accouplements doivent correspondre à l'opération $* = \Lotimes_k$ dans $D_{\geq 0}(k\text{-Mod})$, admettant $k$ comme objet unité, et
\[
\pi_0(K_{\bullet}) \;\simeq\; H_0(K_{\bullet}).
\]
Donc on doit avoir
\[
\Homext(a,b) \Lotimes_k \Homext(b,c) \;\longrightarrow\; \Homext(a,c),
\]
où cette fois Hom ext désigne l'objet de $D_{\geq 0}(k\text{-Mod})$, au lieu d'un objet de Hot seulement, ou de $\Hot_{\bullet}$. On utilise les foncteurs
\[
D_{\geq 0}(k\text{-Mod}) \longrightarrow \Hot_{\bullet} \longrightarrow \Hot
\]
\msp{page 20}
pour faire le lien avec la théorie ponctuée dans $\Hot_{\bullet}$, ou \texttt{[la]} théorie générale dans Hot. Si on désigne par
\[
\gamma :\ D_{\geq 0}(k\text{-Mod}) \;\longrightarrow\; \Hot_{\bullet}
\]
le premier de ces foncteurs, il ne commute pas aux produits $*$ ($\Lotimes_k$ et $\times$ respectivement), mais on a, pour $P_{\bullet}$, $Q_{\bullet}$ variable dans $D_{\geq 0}(k\text{-Mod})$, une flèche fonctorielle canonique % [sic] variable -> variables, agreement with P, Q
\[
\gamma(P_{\bullet} \Lotimes_k Q_{\bullet}) \xleftarrow{\ \mathit{can.}\ } \gamma(P_{\bullet})\times\gamma(Q_{\bullet})
\]
`$k$-bilinéaire', ce qui montre que de donner des accouplements
\[
P_{\bullet} * Q_{\bullet} \xrightarrow{\ c_k\ } R_{\bullet},
\]
d'où
\[
\gamma(P_{\bullet} * Q_{\bullet}) \xrightarrow{\ \gamma(c_k)\ } \gamma(R_{\bullet}), \tag*{$(*)$}
\]
est l'un et l'autre plus fin que de se donner seulement
\[
\gamma(P_{\bullet}) * \gamma(Q_{\bullet}) \xrightarrow{\ c_{\bullet}\ } \gamma(R_{\bullet}), \tag*{$(*')$}
\]
vu qu'une donnée $(*)$ impliquera une donne $(*')$ par % [sic] donne -> donnée, missing letters
\[
\begin{tikzcd}
 & \gamma(P_{\bullet} * Q_{\bullet}) \arrow[dr, "\gamma(c)"] & \\
\gamma(P_{\bullet}) * \gamma(Q_{\bullet}) \arrow[ur, "\text{can.}"] \arrow[rr, "c_{\bullet}"] & & \gamma(R_{\bullet}).
\end{tikzcd}
\]
Ceci vu, on s'attend bien à avoir une théorie \textit{forte} des accouplements
\printedpage{13}
% head: A. Grothendieck : Dérivateurs, Ch. IX, Édition des Universités de Montpellier II et Paris VII
% foot: Version 12 mars 2002
\msp{page 21}
\[
\Homext_k(a,b) \Lotimes_k \Homext_k(b,c) \;\longrightarrow\; \Homext_k(a,c),
\]
d'où on déduira les accouplements plus faibles dans $\Hot_{\bullet}$ (au lieu de $D_{\geq 0}(k\text{-Mod})$)
\[
\Homext_{\Hot_{\bullet}}(a,b) \times \Homext_{\Hot_{\bullet}}(b,c) \;\longrightarrow\; \Homext_{\Hot_{\bullet}}(a,c).
\]

3$^o$) Revenant au cas $\M$ ponctuée sans plus, \textsc{Quillen} définit, pour $a$, $b$ dans $\M$,
\[
\pi_1(a,b) \;=\; \pi_1(a,b\,;0_a,0_b)
\]
comme bifoncteur sur $(\M\Sigma^{-1})^o\times(\M\Sigma^{-1})$ (avec les notations de 1$^o$)). Si ce que j'avance sans preuve dans 1$^o$) est vrai, il doit s'ensuivre qu'on a un isomorphisme canonique
\[
\pi_1(a,b) \;\simeq\; \pi_1(\Ch^{\Sigma}_{\infty}(a,b)),
\]
où $\Ch^{\Sigma}_{\infty}(a,b)$ est vu comme objet de $\Hot_{\bullet}$, donc son $\pi_1$ est défini. On a en tous cas, sans aucune référence à des cofibrations et des fibrations, un bifoncteur à valeurs dans Groupes, composé dans
\[
(\M\Sigma^{-1})^o\times(\M\Sigma^{-1}) \xrightarrow{\ \ (a,b)\longmapsto \Ch^{\Sigma}_{\infty}(a,b)\ \ } \Hot_{\bullet} \xrightarrow{\ \pi_1\ } (\text{Groupes}).
\]
Si pour $a$ fixé, $b$ variable,
\[
b \;\longmapsto\; \pi_1(a,b)
\]
est représentable, on note l'objet qui le représente par $\Sigma(a)$ (\textit{suspension} -- ne pas confondre ce $\Sigma$ avec $\Sigma \subseteq \mathrm{Fl}(\M)$, qu'il vaut mieux maintenant noter $W$, comme `weak equivalence'). Donc

\msp{page 22}

\[
\mathrm{Hom}(\Sigma a,b) \;\simeq\; \pi_1(a,b),
\]
et $\Sigma a$ dépend fonctoriellement de $a$, sous réserve d'existence. S'il existe toujours, on a
\[
\Sigma :\ \M W^{-1} \;\longrightarrow\; \M W^{-1}.
\]
On définit de même, sous réserve d'existence,
\[
\Omega :\ \M W^{-1} \;\longrightarrow\; \M W^{-1}
\]
par
\[
\mathrm{Hom}(a,\Omega b) \;\simeq\; \pi_1(a,b),
\]
et si tous les deux existent, on a la formule d'adjonction
\[
\mathrm{Hom}(\Sigma a,b) \;\simeq\; \mathrm{Hom}(a,\Omega b) \qquad (\simeq \pi_1\Homext_{\Hot_{\bullet}}(a,b)).
\]

% ---- from Der-ChIX-678_p016-018.tex
\printedpage{14}
% head: A. Grothendieck : Dérivateurs, Ch. IX, Édition des Universités de Montpellier II et Paris VII
% foot: Version 12 mars 2002

L'existence de
\[
\begin{aligned}
\Omega :\ & \mathcal{M}W^{-1} \ \longrightarrow\ \mathcal{M}W^{-1}, \quad \text{ou de}\\
\Sigma :\ & \mathcal{M}W^{-1} \ \longrightarrow\ \mathcal{M}W^{-1},
\end{aligned}
\]
voire même des deux à la fois, apparaît comme une hypothèse \cercle{a} extrêmement forte sur $(\mathcal{M}, W)$. Cette condition \textit{plus} celle \cercle{b} que les $\Chu_{\infty}(a,b)$ soient homotopes à des catégories de Thomason-Kan, ou (ce qui revient au même, à peu de choses près?) que $Y \longmapsto \mathrm{Hom}_{\mathcal{M}(Y)W_Y^{-1}}(a_Y, b_Y)$ transforme Hot-équivalences de Cat en bijections (ne serait-ce que pour $Y$ fixé \dots), semble le test crucial pour savoir si `$(\mathcal{M}, \Sigma)$ \texttt{[plutôt} $(\mathcal{M}, W)$\texttt{]} se prête à une théorie homotopique `dans $\mathcal{M}$' '. Je doute que

\msp{page 23}

ces conditions puissent être satisfaites pour $\mathcal{M}$ finie, sauf dans le cas trivial où $\mathcal{M} = \emptyset$, où $\mathcal{M}\Sigma^{-1}$ \texttt{[plutôt} $\mathcal{M}W^{-1}$\texttt{]} est équivalente à la catégorie ponctuelle.

Un troisième test \cercle{c} qui me paraît important (mais peut-être ça résulte du reste), c'est les formules dans $\Hot_\bullet$
\[
\left\{
\begin{aligned}
\Homext_\bullet(\Sigma a, b) \ &\simeq\ \Homext_\bullet(a, \Omega b)\\
&\simeq\ \Omega\Homext_\bullet(a, b),
\end{aligned}
\right.
\tag*{$(*)$}
\]
où le $\Homext_\bullet$ désigne le Hom externe à valeurs dans $\Hot_\bullet$. Il faut le voir séparément (sous condition d'existence soit de $\Omega$, soit de $\Sigma$) comme deux formules `duales' l'une de l'autre,
\[
\left\{
\begin{aligned}
\Homext_\bullet(a, \Omega b) \ &\simeq\ \Omega\Homext_\bullet(a, b)\\
\Homext_\bullet(\Sigma a, b) \ &\simeq\ \Omega\Homext_\bullet(a, b),
\end{aligned}
\right.
\tag*{$(**)$}
\]
isomorphismes de \textit{groupes} dans $\Hot_\bullet$. (\textbf{NB} $\Omega b$ est un objet groupe de $\mathcal{M}W^{-1}$, car il représente un contrafoncteur à valeurs dans Gr, et de même $\Sigma a$ est un cogroupe de $\mathcal{M}W^{-1}$. Le foncteur $y \longmapsto \Homext(a, y)$ doit commuter aux produits finis, et donc transformer objets groupes en itou. Et de même $x \longmapsto \Homext(x, b)$ doit transformer sommes finies en produits finis, donc cogroupes en groupes (\footnote[8]{Le foncteur $\Omega$ doit provenir d'un endomorphisme du dérivateur $\mathbf{D} = \mathbf{D}_{\mathcal{M},W}$, et celui-ci doit être exact à gauche, et dualement pour $\Sigma$.}).)

Je me suis convaincu que les isomorphismes $(**)$ n'ont de chance d'exister que si on a la condition préalable \cercle{b}.

\newpage

\printedpage{15}
% head: A. Grothendieck : Dérivateurs, Ch. IX, Édition des Universités de Montpellier II et Paris VII
% foot: Version 12 mars 2002

\msp{page 24}

\section*{3\quad Essai de construction du foncteur $\Omega$}

Il n'y a sans doute guère de chance de prouver quelque chose sur le foncteur
\[
\Omega :\ \mathcal{M}W^{-1} \longrightarrow \mathcal{M}W^{-1},
\]
sans tenir en mains une \textit{construction} plus ou moins explicite de celui-ci. Je vais essayer, en m'inspirant de \textsc{Quillen}, de dégager des conditions générales sur $(\mathcal{M}, W)$ (plus générales que celles de \textsc{Quillen}) qui assurent l'existence de $\Omega$, et permettent de le construire. L'idée géométrique remonte à \textsc{Serre-Cartan}. Il s'agit, pour un objet $S$ de $\mathcal{M}$, de construire un objet $\mathrm{Ch}(S)$ \textit{dans $\mathcal{M}$,} jouant le rôle d'un `espace de chemins', et qui doit s'inclure dans un diagramme de factorisation de $\mathrm{diag}_S : S \longrightarrow S \times S$ (on suppose que dans $\mathcal{M}$ les $\varprojlim$ finies, donc les produits finis, existent)
\[
\begin{tikzcd}
S \arrow[d, "i_S"] \arrow[dd, bend right=40, "\mathrm{diag}_S"'] & \\
\mathrm{Ch}(S) \arrow[d, "p_S"] & \\
S \times S & e. \arrow[l]
\end{tikzcd}
\]
La propriété essentielle, c'est que $p$ soit une \textit{`fibration'} (dans un sens convenable, à dégager).

\msp{page 25}

À titre accessoire il peut être utile (peut-être) que $i$ soit une `cofibration' (\footnote[9]{et peut-être que $S \longrightarrow e$ soit également une fibration ?}). On définit $\Omega(S)$ comme la fibre de $\mathrm{Ch}(S)$ sur $S{\times}S$ (par rapport à l'unique morphisme $e = e_{\mathcal{M}} \longrightarrow S\times S$, où $e_{\mathcal{M}}$ est l'objet nul de $\mathcal{M}$).
\[
\begin{tikzcd}[column sep=large]
\mathrm{Ch}(S) \arrow[d] & \Omega S \arrow[l] \arrow[d]\\
S \times S & e. \arrow[l]
\end{tikzcd}
\]

\newpage

\printedpage{16}
% head: A. Grothendieck : Dérivateurs, Ch. IX, Édition des Universités de Montpellier II et Paris VII
% foot: Version 12 mars 2002

C'est un objet de $\mathcal{M}$, qu'on regardera surtout comme un objet de $\mathcal{M}W^{-1}$. Il faut s'assurer qu'on a une bijection fonctorielle en $T$ dans $\mathcal{M}W^{-1}$
\[
\underbrace{\pi_1\Chu^W_\infty(T, S)}_{\overset{\text{déf}}{=}\,\pi_1(T,S)} \ \simeq\ \underbrace{\mathrm{Hom}_{\mathcal{M}W^{-1}}(T, \Omega S)}_{\overset{\text{déf}}{=}\,\pi_0\Chu^W_\infty(T,\Omega S)}.
\]
Pour cela, il s'impose de comparer
\[
\Chu^W_\infty(T, \Omega S) \qquad\quad \text{et} \qquad\quad \underbrace{\Omega_{\Hot_\bullet}(\Chu^W_\infty(T, S))}_{\overset{\substack{\text{`par définition'}\\ \text{pour $\Omega$ de $\mathrm{Hot}_\bullet$}}}{=}\ \Chu_\infty(\mathrm{Ch}^W_\infty(T,S);0,0)} \qquad ,
\]
puisqu'il s'agit de construire une bijection entre leurs $\pi_0$.

% ---- from Der-ChX-789_p001-005.tex
% ===================== page 1 (p001) =====================
\printedpage{}
% no page number printed (title page)
% head: (none)

\begin{center}
{\LARGE \textsf{LES DÉRIVATEURS}}

\bigskip

{\large \textbf{ALEXANDRE GROTHENDIECK}}

\vspace{3em}

{\large \textbf{Chapitre X}}

\medskip

{\large \textbf{Comparaison de} Cat \textbf{avec} $\bDelta^{\wedge}$}
\end{center}

\clearpage

% ===================== page 2 (p002) =====================
\printedpage{}
% no page number printed (editorial note)
% head: (none)

Ce texte a été déchiffré et transcrit en \LaTeXe{} par M. Künzer. Il a été édité par M. Künzer, J. Malgoire, et G. Maltsiniotis. La transcription est aussi fidèle que possible au manuscrit. Pour les quelques corrections évidentes, ou rares commentaires des éditeurs, ainsi que pour la numérotation originale des pages du manuscrit, les caractères de machine à écrire \texttt{[typewriter]} entre crochets sont utilisés. Un point d’interrogation entre crochets signifie que l’on n’est pas sûr du mot qui précède.

Cette édition est provisoire. Les remarques, commentaires et corrections sont bienvenus. Envoyer un message à :

\textbf{[e-mail address removed]}

\hfill G. Maltsiniotis\hspace*{3em}

\clearpage

% ===================== page 3 (p003), printed 1 =====================
\printedpage{1}
% head: A. Grothendieck : Dérivateurs, Ch. X, Édition des Universités de Montpellier II et Paris VII
% foot: Version 9 avril 2002

\msp{page 1}

\section*{1\quad Rappel des notions de D-asphéricité, et lien avec le ‘localisateur fondamental $W_{\mathbf{D}}$’}

Pour un dérivateur $\D$, il s’agit des notions suivantes (cf. I) :

\begin{itemize}
\item[(i)] \textit{Objet} $\D$-\textit{asphérique} de Cat (i.e. $Z \longrightarrow e$ est une $\D$-équivalence). NB Je préfère dire ‘$\D$-asphérique’ au lieu de ‘$\D$-contractile’, qui prête parfois à confusion.

\item[\textcircled{ii}] $\D$-équivalence pour une flèche $f : X \longrightarrow Y$ dans Cat.

\item[(iii)] \textit{Foncteur} $f : X \longrightarrow Y$ $\D$-\textit{asphérique} (i.e. les $X/y$ sont $\D$-asphériques) et $\D$-\textit{coasphérique} (les $y\backslash X$ sont $\D$-asphériques). C’est ce que j’avais appelé ‘$\D$-équivalence stricte’ (ou ‘costricte’), ou foncteur homologiquement $\D$-cofinal (ou cohomologiquement $\D$-cofinal), mais la terminologie ‘asphérique’ me semble la plus simple et la meilleure. (Pêchée dans PS, Pursuing Stacks.)

\item[(iv)] \textit{Foncteur} $f : X \longrightarrow Y$ $\D$-\textit{lisse,} $\D$-\textit{propre} (i.e. tel que les fibres des foncteurs induits $X/x \longrightarrow Y/f(x)$, respectivement $x\backslash X \longrightarrow f(x)\backslash Y$, soient $\D$-asphériques).

\item[(v)] $\D$-\textit{fibration} $f : X \longrightarrow Y$ (notion autoduale), i.e. telle que pour $Y'' \xrightarrow{u} Y'$ au dessus de $Y$, si $Y''$, $Y'$ \texttt{[sont]} Hot-asphériques, alors $u_X : X'' \longrightarrow X'$ dans
\[
\begin{tikzcd}
X \arrow[d] & X' \arrow[l] \arrow[d] & X'' \arrow[l] \arrow[d] \\
Y & Y' \arrow[l] & Y'' \arrow[l]
\end{tikzcd}
\]
est une $\D$-équivalence (\footnote{On dira aussi que $f$ est $\D$-\textit{fibrant.} (\textbf{NB} Il n’y a pas de notion ‘$\D$-cofibrant’ à ma connaissance.)}). La $\D$-\textit{fibration} est dite \textit{triviale} (ou \textit{asphérique}) si et seulement si ses fibres sont asphériques, ou encore si et seulement si $f$ est asphérique, ou si et seulement si $f$ est coaspérique. Cela signifie aussi que $f$ est une $\D$-équivalence \textit{universelle,} % [sic] coaspérique -> coasphérique, letter h missing
\msp{page 2}
i.e. que pour tout $Y' \longrightarrow Y$, $X' = X \times_Y Y' \longrightarrow Y'$ est une $\D$-équivalence.

\item[(vi)] On a aussi des variantes affaiblies de la notion de $\D$-fibrations, savoir les ‘$\D$-\textit{fibrations à gauche’,} ou les $f$ ‘\textit{cohomologiquement} $\D$-\textit{fibrants’,} et de même les ‘$\D$-\textit{fibrations à droite’,} i.e. les $f$ ‘\textit{homologiquement} $\D$-\textit{fibrants’} (\footnote{Nouvelle terminologie : fibrations cohomologiques (resp. homologiques) \textit{faibles,} cf. I, pages 111-126 pour une révision corrigée des diverses notions de ‘$\D$-fibrations’.}). La notion ‘gauche’
\printedpage{2}
% head: A. Grothendieck : Dérivateurs, Ch. X, Édition des Universités de Montpellier II et Paris VII
% foot: Version 9 avril 2002
ou cohomologique pour $f : X \longrightarrow Y$ signifie que les foncteurs $X_y \hookrightarrow X/y$ ($y \in \Ob Y$) et $X/y \longrightarrow X/y'$ (pour $(u : y \longrightarrow y') \in \Fl Y$) sont des $\D$-équivalences. Et dualement avec les $X_y \hookrightarrow y\backslash X$, $y'\backslash X \longrightarrow y\backslash X$ pour la notion ‘droite’ ou ‘homologique’. Les $\D$-fibrations ne sont autres que les $\D$-fibrations gauches (cohomologiques) universelles ou les $\D$-fibrations droites (homologiques) universelles.
\end{itemize}

Finalement, toute cette kyrielle de notions se déduit de la seule notion de $\D$-équivalence (ii). Si je désigne par
\[
W_{\D} \subseteq \Fl(\Cat)
\]
l’ensemble de ces $\D$-équivalences, on trouve ce que dans Pursuing Stacks j’ai appelé un ‘localisateur’ (localiseur) dans Cat (\footnote{Mais attention, je suppose ici que le domaine Diag de $\D$ est Cat tout entier!}). Dans Pursuing Stacks c’était une telle partie
\[
W \subseteq \Fl(\Cat),
\]
et non la donnée d’un dérivateur, qui
\msp{page 3}
était pris comme la donnée de départ fondamental pour faire de l’algèbre homotopique en tous genres. J’avais d’ailleurs développé une kyrielle de notions du même genre (i) à (vi), en partant de cette donnée $W$. Il se pourrait que les énoncés pour les notions que j’ai développé dans I, en termes d’un dérivateur, soient valables plus généralement en termes d’un localisateur $W$, dans la mesure du moins où ils gardent un sens. (Donc pour l’instant je mets à part les énoncés proprement cohomologiques ou homologiques, qui font directement appel à $\D$.) Encore qu’il se pourrait que moyennant des conditions à expliciter sur $W$, qui impliquent que $(\Cat, W)$ définissent bel et bien un dérivateur $\D_{\Cat,W} = \D$, on arrive encore à interpréter $W$ comme un $W_{\D}$, de sorte que les énoncés cohomologiques et homologiques gardent un sens.

J’ai envie d’abord de regarder dans Pursuing Stacks les axiomes que j’ai dégagé pour les localisateurs, et vérifier s’ils sont tous vérifiés pour
\msp{page 4}
$W_{\D}$. Ces conditions L 1 à L 5 sont énoncées dans Pursuing Stacks, §64 (page 164).

\textbf{L 1}\enspace Saturation de $W_D$. O.K.

Cela résulte formellement du fait qu’on a un foncteur
\[
\Cat \xrightarrow{\varphi} C \qquad (= \Cat^o)
\]
tel que pour $u \in \Fl(\Cat)$, $u \in W_{\D} \iff \varphi(u)$ isomorphisme. (Prendre pour $\varphi(X)$ la sous-catégorie de $\D(X)$ formée des $\D$-coefficients \textit{constants} sur $X$).

\textbf{L 2}\enspace “$\bDelta^1$ est $W$-asphérique sur $e$”, i.e. pour tout $X$ dans Cat, $X \times \bDelta^1 \longrightarrow X$ est dans $W$.

Cela signifie aussi que $\bDelta^1 \longrightarrow e$ est une $\D$-fibration triviale, ce qui est bien le cas.

\clearpage

% ===================== page 5 (p005), printed 3 =====================
\printedpage{3}
% head: A. Grothendieck : Dérivateurs, Ch. X, Édition des Universités de Montpellier II et Paris VII
% foot: Version 9 avril 2002

\textbf{L 3}\enspace Si $X \in \Ob\Cat$ a un objet final, alors $X \longrightarrow e$ est dans $W$.

En fait, nous exigeons pour les dérivateurs que ceci soit vrai en supposant seulement $X$ asphérique (au sens absolu), et on en conclut même que $X \longrightarrow e$ est $W$-asphérique, i.e. $X \times Y \xrightarrow{\mathrm{pr}_2} Y$ est dans $W$ pour tout $Y$ dans Cat (\footnote{En fait L 2 $\Longrightarrow$ L 3.}).

\textbf{L 4}\enspace Si $f : X \longrightarrow Y$ est tel que $X/y \longrightarrow Y/y$ sont dans $W$, alors $f$ est dans $W$.

C’est OK.

\msp{page 5}
\textbf{L 5}\enspace Si $X$, $Y$ sont Cat-fibrées sur $S$, et $X \longrightarrow Y$ est un $S$-morphisme \textit{cartésien,} induisant sur les fibres des $X_s \longrightarrow Y_s$ qui sont dans $W$, alors $f$ est dans $W$.

Pour $W_{\D}$, c’est la proposition 6, page 71 de I.

Ces axiomes L 1 à L 5 (équivalent à L 1, L 3$'$ ($\bDelta^1 \longrightarrow e$ est dans $W$) et L 5) sont dans la nature d’axiomes de stabilité, assurant que $W$ est assez grand. Il y a une autre condition qui limite $W$ à ne pas être trop grand :

\textbf{L a)}\enspace Si $f \in W$, alors $\pi_0(f)$ est bijectif.

Ceci n’est pas satisfait pour tous les $W_{\D}$ – p. ex. si $\D$ est le dérivateur \textit{nul,} alors $W_{\D} = \Fl\Cat$! Mais pour les dérivateurs utiles en pratique, cela semble toujours satisfait.

\textbf{L b), c)}\enspace Le foncteur
\[
\Cat \longrightarrow \Cat\, W^{-1}
\]
commute aux sommes finies et aux produits finis (\footnote{Ça marche si $(\Cat, W)$ provient d’un triple de Quillen (avec des Q-fibrations, Q-cofibrations, Q-équivalences faibles).}).

À examiner ce que ça signifie en termes concrets sur $W$ – je n’ai pas vérifié que c’est O.K. pour les $W_{\D}$.

\textbf{L d)}\enspace Si $f : X \longrightarrow Y$ dans $W$, alors $f^o : X^o \longrightarrow Y^o$ aussi.

C’est le corollaire 2, page 75 de I.
\end{document}
